\section{Introduction}
\label{sec:intro}

Type inference with subtyping expresses the typing requirements of a program
as \emph{subtype constraints} $t\le t'$ between type expressions with unknowns
\citep{Mitchell91,FuhMishra90,AikenWimmers93,EifrigSmithTrifonov95,Pottier98}.
An inferred \emph{constrained type} consists of a type and constraints on its
variables. Such constraint sets grow quickly, so simplification removes
redundant constraints \citep{Pottier96,Pottier01,TrifonovSmith96,Rehof97}.
For example, suppose a constrained type includes
$u\le a\to b$, $a\to b\le v$, and $u\le v$.
The first two inequalities imply the third by transitivity, so a simplifier
can remove $u\le v$ without changing the solutions. They do not imply
$v\le u$: the assignment $u=\Bot$, $v=\Top$, and $a=b=\Bot$ satisfies
the first two inequalities and refutes the reversed query.

The general question is \emph{entailment}: given a conjunction $\phi$ and
an inequality $x\le y$, does every solution of $\phi$ satisfy $x\le y$?
We write $\phi\entails x\le y$. Entailment is a universally quantified
implication. It differs from satisfiability, which asks for only one solution.
For structural subtyping, which relates trees of the same shape, entailment
is coNP-complete over finite trees and PSPACE-complete over possibly infinite
trees \citep{HengleinRehof97,HengleinRehof98}. Bottom and top make subtyping
non-structural: $\Bot\le t\le\Top$ can relate trees of different shapes.

TLCA Problem~16, submitted by Jakob Rehof in 1996, asks whether entailment
of subtype constraints over these non-structurally ordered trees is
decidable. Its simple type expressions are built from bottom, top, and a
binary constructor \citep{Rehof96TLCA}. The statement does not specify
variance; the literature treats the constructor as covariant
\citep{NiehrenPriesnitz03}, hereafter NP03. That case is PSPACE-hard
\citep{HengleinRehof98} (\cref{thm:cx-hard} gives another proof) and has
been described as open in subsequent work
\citep{NiehrenPriesnitz03,DudenhefnerMartensRehof17,JiangCuiOliveira25}.
We decide non-structural subtype entailment (NSSE) for every variance of
one constructor.

\begin{theorem}[Main theorem]\label{thm:main}
For every finite arity $n$ and variance map $c:\Dir\to\{0,1\}$, where
$\Dir=\{0,\ldots,n-1\}$,
non-structural subtype entailment over $\{\Bot,f,\Top\}$, where $f$ is the
single constructor of arity $n$, is decidable over arbitrary trees, regular
trees, and finite trees. The decision procedures are uniform in $n,c$.
\end{theorem}

The procedures take flat literals
$z\le f(a_i)_i$, $f(a_i)_i\le z$, $z=\Bot$, and $z=\Top$, with variables
$z,a_i$, and a variable query $x\le y$.
\Cref{prop:flattening} gives the reduction from general finite systems
at positive arity and the direct enumeration at arity zero. The mechanized
procedures take flat systems; flattening is proved here but not mechanized.

The theorem thus answers TLCA Problem~16 for every choice of binary
variance and for each of the three tree domains. Binary covariance is
$c(0)=c(1)=0$; arrows have $c(0)=1,c(1)=0$, with a contravariant domain
and a covariant codomain. Regular entailment coincides with arbitrary-tree
entailment. Finite entailment also coincides with it whenever a finite
solution exists; otherwise finite entailment is vacuous.
A \emph{countermodel} is a solution that refutes the query.
\Cref{fig:overview} summarizes the proof.

\begin{figure}[tbp]
\centering
\small
\hyphenpenalty=10000\relax
\begin{tikzpicture}[
    node distance=5mm,
    >=Stealth,
    redbox/.style={draw, rounded corners=2pt, align=center, text width=0.88\linewidth,
                 inner sep=4pt, minimum height=7mm, font=\small},
    domainbox/.style={draw, dashed, rounded corners=2pt, align=flush left,
                 text width=0.285\linewidth, inner sep=3pt,
                 font=\footnotesize}]
  \node[redbox] (ent) {Every solution of $\phi$ satisfies \mbox{$x\le y$}.};
  \node[redbox, below=of ent] (unsafe)
    {No finite path witnesses a failed inequality\\
     (track order reversals with two copies of each variable; \cref{sec:unsafe,sec:signed}).};
  \node[redbox, below=of unsafe] (spine)
    {Every attempt to force such a path adds inconsistent constraints\\
     (four chains of fresh variables along the path; \cref{thm:four-spine}).};
  \node[redbox, below=of spine] (hole)
    {Neither finite reader has a rejected word without an admitted suffix comparison\\
     (such a word is called a hole; \cref{thm:holes}).};
  \node[redbox, below=of hole] (short)
    {Neither reader has a hole of length at most the elementary bound $H$\\
     (bound the tail, the depth, the number of distinct blocks, and the length; \cref{sec:tail,sec:depth,sec:search}).};
  \draw[<->] (ent) -- (unsafe);
  \draw[<->] (unsafe) -- (spine);
  \draw[<->] (spine) -- (hole);
  \draw[<->] (hole) -- (short);
  \node[domainbox, below=6mm of short.south west, anchor=north west] (unsat)
    {\emph{No solution}: a finite consistency test detects this case;
     entailment is vacuous.};
  \node[domainbox, right=2mm of unsat] (reg)
    {\emph{Regular trees}: every failure has a countermodel represented by a finite graph.};
  \node[domainbox, right=2mm of reg] (fin)
    {\emph{Finite trees}: test finite satisfiability, then use the same entailment answer.};
\end{tikzpicture}
\caption{The reduction for a satisfiable system $\phi$. Each double arrow is
an equivalence. The dashed boxes explain the satisfiability test and the
transfers to regular and finite trees.}
\Description{Five boxes link entailment to the absence of a finite counterexample path,
the inconsistency of constraints forcing that path, the absence of reader holes,
and a bounded search. Three dashed boxes cover unsatisfiable systems and the
two restricted tree domains.}
\label{fig:overview}
\end{figure}

\paragraph{Why Entailment Is Hard and What Makes It Decidable}
A failed tree inequality can be witnessed by a finite path, written as a
word over argument positions. After variance is accounted for, a word $w$
is \emph{left unsafe} if, in some solution, $x$ has a top leaf on a prefix
of $w$ but $y$ has none, and \emph{right unsafe} if $y$ has a bottom leaf on
a prefix of $w$ but $x$ has none. Entailment fails exactly when some word
is left or right unsafe.

Suffix comparisons arise when constraints connect subtrees at different
cuts of this path. In the covariant case, let $s<e$ be cuts of $w$, and
suppose that some input variable lies above the subtree of $x$ at cut $e$
and below the subtree of $y$ at cut $s$. Every top leaf of $x$ on $w$ then
lies at or beyond cut $e$, so the subtree of $x$ at cut $e$ has a top leaf
on a prefix of the remaining suffix $w[e:]$. The variable passes such a
leaf up to the subtree of $y$ at cut $s$, again on a prefix of $w[e:]$. If
$w[e:]$ is a prefix of $w[s:]$, this leaf of $y$ lies on $w$, so $w$ is not
left unsafe after all. A pair of cuts $s<e$ with $w[e:]$ a prefix of
$w[s:]$ is a \emph{comparison} of $w$; equivalently, $w[s:]$ has period
$e-s$. Under variance, comparing a subtree of $x$ with a subtree of its own
signed copy gives an analogous obstruction.

For example, take a binary covariant constructor, the single literal
$f(y,x)\le y$, and the query $x\le y$. In every solution $y\ne\Bot$, and
whenever $y$ is a constructor, it lies below its own $0$-child and $x$ lies
below its $1$-child. Hence, if $y$ has a bottom leaf on a word $w$, then
$w=0^i1z$, the nodes of $y$ along $0^i$ are constructors, and the subtree
of $y$ at cut $i+1$ lies above $x$. This is the order dual of the situation
above, with cuts $0<i+1$ and $x$ itself as the connecting variable: the
bottom leaf of $y$ lies in that subtree on a prefix of $z$, and it forces a
bottom leaf of $x$ on a prefix of $z$. If $(0,i+1)$ is a comparison, that is, if
$z$ is a prefix of $w$, then $x$ also has a bottom leaf on $w$, and $w$ is
not right unsafe. Conversely, a word $0^i1z$ for which $(0,i+1)$ is not a
comparison is right unsafe (\cref{sec:general-holes}). The right-unsafe
words must therefore avoid a periodic comparison between their suffixes,
and they do not form a regular set: those in $0^*10^*$ are exactly the
words $0^i10^j$ with $j>i$.

A finite \emph{reader} records the constraints along a word. It
\emph{rejects} a word when none of the tests at single cuts rules the word
out, so rejection marks a candidate counterexample word. A comparison is
\emph{admitted} when its two reader states certify an obstruction of the
kind just described. A \emph{hole} is a rejected word with no admitted
comparison. For a satisfiable system, the holes of two such readers are
exactly the left-unsafe and the right-unsafe words. Ordinary loop deletion
can preserve the final reader state while creating a new admitted
comparison (\cref{ex:loop-deletion}). Preserving the run is therefore
insufficient.

Readers constructed from subtype constraints satisfy a
\emph{rejected-path property}: the rejected cuts of a word form an
interval, and every admitted comparison ends no later than the first
rejected cut. The reason is semantic. In the left label reader, at a rejected
cut, no variable that the constraints place above the subtree of $x$ is
bounded above by a constructor. From the next cut on, the constraints
therefore bound the subtree of $x$ by no variable at all, whereas every
admission needs such a bound at its later cut (\cref{sec:interface}). The
right label reader is dual. \Cref{tab:reader-hypotheses} records which bounds
need this property.

\paragraph{From Constraints to a Finite Search}
To handle variance, normalization exchanges $\Bot$ and $\Top$ at paths
of odd polarity, that is, paths with an odd number of contravariant edges.
Positive and negative copies of each variable then give a covariant
system; the solutions that matter assign dual trees to the two copies. We
obtain such a solution by selecting labels on a constructor shape common to both
copies. Four chains of fresh constraints along the candidate word, called
\emph{spines}, impose the two unsafety conditions and their signed
counterparts. Classifying the inconsistencies of this extension gives the
reader. For $k$ input variables, a state stores two sets of the $2k$ signed
variables and one bit recording an earlier obstruction, so there are
$N=2\cdot4^{2k}=2^{4k+1}$ states.

The finite search uses three bounds.
\begin{enumerate}[leftmargin=*]
\item \emph{Tail} (\cref{sec:tail}). If a reader with the rejected-path
  property has a hole, then it has one whose \emph{rejected tail}, the part
  after the first rejected cut, has length linear in $N$. If a cycle of
  rejected states allows a second letter, repeating the cycle and then
  taking the other letter breaks every relevant period; otherwise
  continuation within rejected states is forced, and a long periodic tail
  can be truncated.
\item \emph{Depth} (\cref{sec:depth}). For every reader, a suffix
  comparison of a rejected word that joins equal states has gap at most
  that word's rejected-tail length. Cutting a word at the occurrences of
  its last letter, and repeating this on the resulting word of blocks,
  gives a hierarchy of block decompositions. The gap bound limits its
  depth, quadratically in $N$ when the tail is linear in $N$.
\item \emph{Length} (\cref{sec:search}). For every reader, a hole of
  bounded depth can be shortened. Merging interchangeable letters bounds
  the number of distinct blocks at each positive level, and the blocks of
  a level are then shortened simultaneously. Finite summaries preserve
  transitions, admissions, and block equalities.
\end{enumerate}
Together, the three bounds give a hole whose length is triply exponential
in $k$ whenever some hole exists. For arbitrary readers, the remaining
question is whether a computable rejected-tail bound exists
(\cref{sec:general-holes}).

\paragraph{Contributions and Scope}
\begin{itemize}[leftmargin=*]
\item A decision procedure for one constructor of any finite arity and
  variance over arbitrary, regular, and finite trees (\cref{thm:main}).
\item A reduction of non-entailment to holes of finite readers
  (\cref{sec:variance,sec:spines,sec:holes}). Closure-based satisfiability
  characterizations are classical
  \citep{KozenPalsbergSchwartzbach94,PalsbergOKeefe95,PalsbergWandOKeefe97}.
  The new semantic ingredients are a least constructor shape common to the
  two signed copies of each variable, the selection of labels on that shape
  to handle variance, and the four-spine analysis that defines the reader.
\item Alphabet-generic bounds on tails, depth, and length
  (\cref{sec:tail,sec:depth,sec:search}); the latter two apply to every
  reader. For readers with the rejected-path property, they bound the
  length of a shortest hole doubly exponentially in a polynomial of $N$,
  hence triply exponentially in $k$ for the label readers
  (\cref{thm:hole-decision}).
\item Countermodel transfers (\cref{sec:decision,app:domains}), including
  coincidence of finite and arbitrary-tree entailment whenever a finite
  solution exists. That solution excludes cycles in the base constraints;
  the spines preserve this exclusion and permit a finite countermodel.
\item Complexity bounds in each domain (\cref{sec:complexity}). A search
  that stores one level of the block hierarchy at a time places entailment in
  2-EXPSPACE and 3-EXPTIME (\cref{thm:complexity}); constructors of arity at
  most one admit stronger bounds (\cref{app:small-arity}). For a covariant
  constructor of every fixed arity
  $n\ge2$, entailment is PSPACE-hard (\cref{thm:cx-hard}), which for $n=2$
  reproves the non-structural hardness theorems of \citet{HengleinRehof98},
  and shortest holes can have length $2^{\Omega(k)}$ (\cref{prop:cx-length}).
\item An end-to-end Lean~4 mechanization, for flat constraints, of the three
  procedures and their correctness, and of the combinatorial content of the
  complexity bounds: the summary search with its configuration bound, the
  one-letter bound, the hardness reduction, and the exponentially long holes
  (\cref{sec:mechanization,app:lean}).
\item A counterexample to Proposition~16 of \citet{NiehrenPriesnitz03}
  (\cref{app:np03}); \cref{sec:errata} explains which published proofs
  depend on that proposition.
\end{itemize}
The proof does not use the cap-automaton construction or its
back-translation from \citet{NiehrenPriesnitz03}. An earlier version of this
work proved the binary covariant case through a mechanized form of the
forward cap-automaton characterization; \cref{sec:experience} describes that
proof and explains why we replaced it.

\paragraph{How to Read This Paper}
\Cref{sec:background} fixes the semantics and the input language.
\Cref{sec:variance,sec:spines,sec:holes} turn subtype constraints into a
word problem. They build on one another and end with the label readers and
their rejected-path property. \Cref{sec:tail,sec:depth,sec:search} solve
that word problem; their principal results are the tail, depth, and
bounded-length theorems (\cref{thm:brt,thm:depth,thm:search}). From the
earlier sections they use only the definitions of comparisons, readers, and
holes and the rejected-path property, all in \cref{sec:holes}, so a reader
interested mainly in the word combinatorics can start there.
\Cref{sec:decision} assembles the decision procedures and the domain
transfers. \Cref{sec:complexity} replaces the explicit enumeration of words
by a search over summaries; it uses the lettered monitors and summaries of
\cref{sec:search}, the reader construction of \cref{sec:holes}, and the
statements of \cref{sec:decision}. \Cref{sec:mechanization} describes the
Lean artifact, \cref{sec:related} discusses related work,
\cref{sec:experience} compares the present proof with the earlier binary
proof, and \cref{sec:future} discusses open problems. The appendices supply
the spine proofs, the preservation arguments for cores and compression, the
domain transfers, the resource accounting and lower-bound details for the
complexity results, and the correspondence with the Lean development; they
end with a navigation table of proof distinctions and the counterexample to
NP03's Proposition~16.
